% ID: FIS-FLU-PID-002
% Titolo: Legge della pressione idrostatica
% Disciplina: Fisica
% Area: Fluidostatica
% Argomento: Pressione idrostatica
% Sottoargomento: Relazione tra pressione e profondita
% Classe: Terza liceo artistico
% Difficolta: 2
% Tipo: quesito_teorico
% Risultato: p = p_0 + rho g h; pressione relativa Delta p = rho g h
% Tempo_stimato: 5 min
% Competenze: richiamo di una legge fisica, uso del linguaggio scientifico, interpretazione delle grandezze
% Prerequisiti: pressione, densita, profondita in un liquido
% Tag: fisica, fluidostatica, pressione_idrostatica, profondita, liquidi
% Fonte: verifica_fisica_fluidostatica_anonimizzata
% Autore: Riccardo Carli
% Licenza: CC-BY-SA-4.0

\begin{esercizio}
Scrivi l'espressione della legge fisica che lega la pressione di un liquido
alla profondita in cui viene misurata. Specifica il significato delle grandezze
che compaiono nella formula.
\end{esercizio}

\begin{soluzione}
La legge della pressione idrostatica afferma che, in un liquido fermo, la
pressione aumenta linearmente con la profondita:
\[
p=p_0+\rho gh.
\]

In questa formula:
\[
\begin{aligned}
p &\text{ e la pressione alla profondita } h,\\
p_0 &\text{ e la pressione sulla superficie libera,}\\
\rho &\text{ e la densita del liquido,}\\
g &\text{ e l'accelerazione di gravita,}\\
h &\text{ e la profondita misurata dalla superficie libera.}
\end{aligned}
\]

Se si considera solo l'aumento di pressione dovuto al liquido, cioe la
pressione relativa rispetto alla superficie, si scrive:
\[
\Delta p=\rho gh.
\]
\end{soluzione}
